\subsection{The Day's Prayer}\label{sec:prayer}

To watch is to pray, and the Church has a prayer for every part of the day. The
Liturgy of the Hours is not the priest's book alone. The Second Vatican Council
wanted the principal Hours, and Vespers first of them, prayed together in the
churches on Sundays and the greater feasts by the care of their pastors, and it
urged lay people to take up the Office as well, with their priests, with one
another, or by themselves (\emph{Sacrosanctum Concilium}, n.~100). When Pope
St Paul VI gave the reformed Office to the Church on 1 November 1970, in his
constitution \emph{Laudis canticum}, he called it the prayer of the whole people
of God, named Morning and Evening Prayer as its two hinges, and offered the
Hours to all the faithful, even those not bound to say them.

In Advent each Hour has its own colour. The Office of Readings reads a passage
of Scripture and a second reading, most often from a Father of the Church, but
also from saints and teachers of later ages; on Sundays, solemnities and feasts
it ends with the \latin{Te Deum}, which is not said on the weekdays of the
season. Morning Prayer
sings every day the Canticle of Zachary, the song of the father of John the
Baptist over his newborn son. Evening Prayer sings the Canticle of Mary, and
from 17 to 23 December its antiphon is one of the great O antiphons, one for
each evening, which the Roman Breviary sang at Vespers, its Evening Prayer, in
those days before the reform as well. Night Prayer closes the day with the Canticle of
Simeon. Whoever can say Morning and Evening
Prayer in Advent prays at dawn and dusk with the whole Church; whoever can say
only one Hour should choose one of those two.

The Fathers read in the Office of Readings are the best company the season
offers, and some of them teach prayer itself. On the Friday of the first week
the Church reads the opening of St Anselm's \emph{Proslogion}, a book that
begins by calling the reader away from his business:

\begin{quote}
Up now, slight man! flee, for a little while, thy occupations; hide thyself, for
a time, from thy disturbing thoughts. Cast aside, now, thy burdensome cares, and
put away thy toilsome business. Yield room for some little time to God; and rest
for a little time in him. Enter the inner chamber of thy mind; shut out all
thoughts save that of God, and such as can aid thee in seeking him; close thy
door and seek him. \ldots{} Teach me to seek thee, and reveal thyself to me,
when I seek thee, for I cannot seek thee, except thou teach me, nor find thee,
except thou reveal thyself. Let me seek thee in longing, let me long for thee in
seeking; let me find thee in love, and love thee in finding.\par
{\smallskip\noindent\small ---St Anselm, \emph{Proslogion}, chapter 1\par}
\end{quote}

\noindent Anselm's ``little while'' is the measure of most lay prayer. A few
minutes, set apart and guarded, are enough to begin.

For the hours that cannot be set apart St Augustine has a word. On the Friday of
the third week, in the years when it falls before 17 December, the Office reads
his exposition of the thirty-seventh psalm (the thirty-eighth in the
Hebrew numbering), at the verse ``Lord, all my desire is before thee, and
my groaning is not hidden from thee'':

\begin{quote}
For it is thy heart's desire that is thy prayer; and if thy desire continues
uninterrupted, thy prayer continueth also. For not without a meaning did the
Apostle say, ``Pray without ceasing.''\par
{\smallskip\noindent\small ---St Augustine, on Psalm 37\par}
\end{quote}

Then there is the Bible on the table. The Lectionary's introduction explains
the plan of the weekday readings: until 16 December the prophet Isaiah is read
in the order of his book, with Gospels chosen to match him; from the Thursday of
the second week the Gospels turn to John the Baptist; and from 17 December the
first chapters of St Matthew and St Luke tell the events that led straight to
the Lord's birth (nn.~93--94). Isaiah is the prophet of Advent. A household that
reads the day's first reading together, slowly, every evening, will have heard by
Christmas many of the great prophecies of the Messiah, in the order the Church
sets them.

The simplest of all the season's prayers is the Angelus, said at morning, noon
and evening in honour of the Incarnation. Paul VI, in his apostolic exhortation
on devotion to Mary, urged that it be kept, said that it needs no revision, and
noted that it leads the one who prays it from the angel's message to the Passion
and the Resurrection (\emph{Marialis cultus}, n.~41). A family that says the
Angelus through Advent will pray the words the priest prays at the altar as the
Collect of the Fourth Sunday (Section~\ref{sec:fourth-sunday}). The Angelus is
given in Latin and English in Appendix~\ref{sec:prayers}.

In the last week the Church's prayer gathers itself. The \emph{Directory on
Popular Piety and the Liturgy} observes that the old Christmas novena grew up as
a way of opening the treasures of the Church's worship to people who found it
hard to enter, and it proposes that Evening Prayer from 17 to 23 December, with
the great antiphons, be kept in parishes with more solemnity and the people
invited to it: that, it says, would be a true Christmas novena (n.~103). Where a parish sings Evening Prayer on those evenings, go.
